% ID: SEL-D0047-OB-E04
% ESTADO: cotejada. Segunda pasada visual y borrador compilado revisados.
% Fuente: Selectividad/sel_2014_matematicas_aplicadas/Matematicas_CCSS_SEPT.pdf
% SHA-256: ee5bd8e23465d55ed58339e7d9adf3d58423ac2966b53d92e39a2f9b35558ab1
% Página: 2; opción B; ejercicio original 4.

El peso de los huevos de una granja sigue una ley Normal de media desconocida y desviación típica 1.23 gramos. Para estimar la media poblacional se ha tomado una muestra de dos docenas de huevos que han dado un peso total de 1615.2 gramos.
\begin{enumerate}[label=\alph*)]
\item \textbf{(1.75 puntos)} Halle un intervalo de confianza, al 96\%, para la media poblacional.
\item \textbf{(0.75 puntos)} Con el mismo nivel de confianza anterior, si nos exigieran que el intervalo tuviera una amplitud máxima de 0.8, \textquestiondown{}de qué tamaño, como mínimo, habría que tomar la muestra?
\end{enumerate}
